\documentclass[11pt]{amsart}

\usepackage[T1]{fontenc}
\usepackage{lmodern}
\usepackage{microtype}
\usepackage{amsmath,amssymb,amsthm}
\usepackage[hidelinks]{hyperref}

\hypersetup{
  pdftitle={The Polygonal Bigraphs P_{8s,4s-2} Are 2-Edge-Hamiltonian},
  pdfsubject={Hamilton cycles in polygonal bigraphs},
  pdfkeywords={Hamilton cycle, polygonal bigraph, 2-edge-Hamiltonian}
}

\newtheorem{theorem}{Theorem}
\newtheorem{corollary}[theorem]{Corollary}
\newtheorem{proposition}[theorem]{Proposition}

\newcommand{\Z}{\mathbb{Z}}

\title[The polygonal bigraphs $P_{8s,4s-2}$]
{The Polygonal Bigraphs $P_{8s,4s-2}$ Are 2-Edge-Hamiltonian}

\date{}

\subjclass[2020]{05C45}
\keywords{Hamilton cycle, polygonal bigraph, 2-edge-Hamiltonian}

\begin{document}

\begin{abstract}
Simmons conjectured that every polygonal bigraph $P_{m,k}$ is
2-edge-Hamiltonian. We prove the conjecture for the infinite family
$P_{8s,4s-2}$, $s\geq 1$, by an explicit four-edge switch and cyclic
rotations; the cases not already covered by Simmons are exactly
$s\equiv 1\pmod 3$ with $s\geq 4$. For those
parameters the graphs lie in the arithmetic residue
$2<\gcd(m+1,k+1)<k+1$ isolated in \cite[p.~109]{Simmons2020}. The first
such case is $P_{32,14}$; the table of the addendum \cite{Simmons2020a}
leaves that cell blank, so our theorem supplies the one case needed for
that addendum's assertion that all $P_{m,k}$ on $66$ or fewer vertices
are 2-edge-Hamiltonian. The remaining parameters are not addressed.
\end{abstract}

\maketitle

\section{The construction}

A graph is \emph{2-edge-Hamiltonian} if every pair of distinct edges lies
on a common Hamilton cycle. Simmons defined $P_{m,k}$ from a cycle on
$2m$ vertices by joining each odd vertex, in a one-based clockwise
labeling, to the even vertex $2k+1$ positions ahead
\cite[pp.~101--102]{Simmons2020}. He conjectured that every $P_{m,k}$ is
2-edge-Hamiltonian \cite[p.~101]{Simmons2020}. The map $v\mapsto
-v\pmod{2m}$ is an isomorphism between Simmons's orientation and the one
used below, in which each even vertex is joined to the odd vertex $2k+1$
positions ahead, so we may work with the latter.

\begin{theorem}\label{thm:main}
For every integer $s\geq 1$, the polygonal bigraph $P_{8s,4s-2}$ is
2-edge-Hamiltonian.
\end{theorem}

\begin{proof}
Put
\[
  m=8s,
  \qquad
  k=4s-2=\frac m2-2.
\]
Work modulo $2m$ on vertices and modulo $m$ on subscripts. The three
natural perfect matchings of $P_{m,k}$ are
\[
  A_i=\{2i,2i+1\},
  \qquad
  B_i=\{2i+1,2i+2\},
  \qquad
  C_i=\{2i,2i+2k+1\}.
\]
Write $A=\{A_i:i\in\Z/m\Z\}$, and similarly for $B$ and $C$.

The outer cycle $A\cup B$ is Hamilton. The subgraph $B\cup C$ is also
Hamilton, because a chord followed by a $B$-edge sends the even subscript
\[
  i\longmapsto i+k+1,
\]
and
\[
  \gcd(k+1,m)=\gcd(4s-1,8s)=1.
\]
These two Hamilton cycles cover every edge pair except pairs of the form
$\{C_p,A_q\}$.

A chord followed by an $A$-edge sends the even subscript $i$ to $i+k$.
Since
\[
  \gcd(k,m)=\gcd(4s-2,8s)=2,
\]
the subgraph $A\cup C$ consists of two cycles, one on the even subscripts
and one on the odd subscripts. Define
\[
\begin{split}
  H={}&\bigl(A\setminus\{A_k,A_{k+m/2}\}\bigr)
       \cup\bigl(C\setminus\{C_{k+1},C_{k+1+m/2}\}\bigr)\\
     &\cup\{B_{k-1},B_k,B_{k-1+m/2},B_{k+m/2}\}.
\end{split}
\]
We show that $H$ is a Hamilton cycle. For $\varepsilon\in\{0,1\}$,
put $i_t=\varepsilon+tk$ in $\Z/m\Z$. The corresponding component of
$A\cup C$ has cyclic order
\[
  (2i_0,2i_1+1,2i_1,2i_2+1,\ldots,
   2i_{m/2-1},2i_0+1).
\]
Since
\[
  \frac m4 k\equiv \frac m2\pmod m,
\]
the two deleted edges in each component occur halfway around that
component. Reading the two cyclic orders gives four spanning paths with
endpoint pairs
\[
\begin{array}{ll}
  \{2k+1,2k+m\}, & \{2k,2k+m+1\},\\[2pt]
  \{2k+2,2k-1\}, & \{2k+m-1,2k+m+2\}.
\end{array}
\]
Here the second line uses $k=m/2-2$, so
\[
  4k+3=2k+m-1,
  \qquad
  4k+m+3\equiv 2k-1\pmod{2m}.
\]
The four added edges join the paths in the cyclic order
\[
\begin{split}
  2k+1
  &\xrightarrow{B_k}2k+2
  \leadsto 2k-1
  \xrightarrow{B_{k-1}}2k
  \leadsto 2k+m+1\\
  &\xrightarrow{B_{k+m/2}}2k+m+2
  \leadsto 2k+m-1
  \xrightarrow{B_{k-1+m/2}}2k+m
  \leadsto 2k+1.
\end{split}
\]
Thus $H$ is spanning, connected, and $2$-regular, hence Hamiltonian.

Now let $C_p$ and $A_q$ be arbitrary and put $d=q-p\in\Z/m\Z$. Choose
\[
  i\notin
  \{k+1,k+1+m/2,k-d,k+m/2-d\}.
\]
Such an $i$ exists because $m\geq 8$. Then $H$ contains both $C_i$ and
$A_{i+d}$. The rotation $\rho(v)=v+2$ is an automorphism satisfying
\[
  \rho^r(C_i)=C_{i+r},
  \qquad
  \rho^r(A_j)=A_{j+r}.
\]
Taking $r=p-i$, the Hamilton cycle $\rho^r(H)$ contains $C_p$ and
$A_{p+d}=A_q$. Hence every remaining edge pair lies on a Hamilton cycle.
\end{proof}

\begin{corollary}\label{cor:residue}
For every integer $t\geq 1$, the graph $P_{24t+8,12t+2}$ is
2-edge-Hamiltonian and satisfies
\[
  2<\gcd(m+1,k+1)<k+1.
\]
In particular, the first member is $P_{32,14}$.
\end{corollary}

\begin{proof}
Set $s=3t+1$ in Theorem~\ref{thm:main}. Then
\[
  m=24t+8,
  \qquad
  k=12t+2,
  \qquad
  \gcd(m+1,k+1)=3.
\]
Since $k+1=12t+3>3$, the strict inequalities follow. This is precisely
the residue $2<\gcd(m+1,k+1)<k+1$ left open in \cite[p.~109]{Simmons2020}
and re-examined in \cite[p.~102]{Simmons2020a}.
\end{proof}

\begin{proposition}\label{prop:noniso}
For every $s\geq 1$, the graph $P_{8s,4s-2}$ is not isomorphic to
$P_{8s,k'}$ for any $k'\neq 4s-2$ in the canonical range
$1\leq k'\leq 4s-1$.
\end{proposition}

\begin{proof}
For fixed $m$, the bipartite adjacency block of $P_{m,k}$ is a
weight-three circulant with support affinely equivalent to
\[
  S_k=\{0,1,k+1\}\subseteq \Z/m\Z.
\]
An isomorphism between two connected bipartite graphs either preserves or
interchanges their parts; interchanging the parts transposes the block and
replaces its support by its negative. It follows from
\cite[Theorem~1.1]{WiedemannZieve} that isomorphic polygonal bigraphs have
affinely equivalent supports.

For the family above, put
\[
  a=k+1=\frac m2-1.
\]
Then $a$ is a unit modulo $m$, $a-1$ is not, and
\[
  a^2\equiv 1\pmod m.
\]
Consequently, after normalizing an ordered pair at unit distance to
$0,1$, the third element is one of
\[
  a,\qquad 1-a,\qquad a^{-1},\qquad 1-a^{-1}.
\]
Since $a^{-1}=a$, only $a$ and $1-a$ occur. If
$S_{k'}=\{0,1,b\}$ with $b=k'+1$ and
$2\leq b\leq m/2$ is affinely equivalent to $S_k$, then
\[
  b\in\{a,1-a\}.
\]
But $1-a\equiv m/2+2\pmod m$, which is outside this range. Hence
$b=a$ and $k'=k$.
\end{proof}

The addendum concludes that ``all $P_{m,k}$ on 66 or fewer vertices are
2-edge-Hamiltonian'' \cite[p.~103]{Simmons2020a}. Its table
carries no entry at $(m,k)=(32,14)$, although $2<\gcd(33,15)=3<15$; by
Proposition~\ref{prop:noniso} the graph $P_{32,14}$ is alone in its
isomorphism class, consistently with the count of $9$ classes for $m=32$
in \cite[Table~3]{Simmons2020}, and
its $64$ vertices put it beyond the $56$ or fewer
vertices reported in \cite[p.~115]{Simmons2020}.
Theorem~\ref{thm:main} therefore supplies the one case needed for that
66-vertex assertion. The next pair of that residue in our family,
$P_{56,26}$, lies beyond the addendum's range. The companion note
\cite{Companion} in this collection gives an independent covering
certificate for $P_{32,14}$ and settles the isomorphism class
$P_{34,9}\cong P_{34,14}$, which the addendum records as not known to be
2-edge-Hamiltonian.

\begin{thebibliography}{9}

\bibitem{Simmons2020}
G.~J. Simmons,
\emph{Two related (?) 2-edge-Hamiltonian bigraph conjectures},
Bull. Inst. Combin. Appl. \textbf{88} (2020), 98--117.

\bibitem{Simmons2020a}
G.~J. Simmons,
\emph{Two related (?) 2-edge-Hamiltonian bigraph conjectures addendum},
Bull. Inst. Combin. Appl. \textbf{89} (2020), 102--103.

\bibitem{Companion}
\emph{Two residual cases of Simmons's polygonal bigraph conjecture},
companion note in this collection, file
\texttt{residual-polygonal-bigraphs-\allowbreak close-simmons-finite-range}.

\bibitem{WiedemannZieve}
D.~Wiedemann and M.~E. Zieve,
\emph{Equivalence of sparse circulants: the bipartite \'Ad\'am problem},
arXiv:0706.1567, 2007.

\end{thebibliography}

\end{document}